\documentclass[11pt,a4paper]{article}
\usepackage[margin=1in]{geometry}
\usepackage{amsmath,amssymb,amsfonts}
\usepackage{mathtools}
\usepackage{lmodern}
\usepackage[T1]{fontenc}
\usepackage[utf8]{inputenc}
\usepackage{textcomp}
\usepackage{microtype}
\usepackage{parskip}
\usepackage{enumitem}
\usepackage{array}
\usepackage{booktabs}
\usepackage{hyperref}
\hypersetup{hidelinks,
  pdftitle={Paper 062 -- Cross-Domain Echo with BH Gamma_cross Collapse},
  pdfauthor={Allen Proxmire}}
\setlength{\parindent}{0pt}
\setlength{\parskip}{6pt plus 2pt minus 1pt}

\title{\vspace{-1cm}\Large\bfseries Paper 062 --- Cross-Domain Echo with BH $\Gamma_\mathrm{cross}$ Collapse}
\author{Allen Proxmire}
\date{May 2026}

\begin{document}
\maketitle

\section*{Preamble --- What This Paper Does NOT Claim}

\begin{enumerate}[noitemsep]
\item This paper does \textbf{not} claim that Q-COMPUTE platforms physically realize black-hole conditions; the echo is at substrate-mechanism level.
\item It does \textbf{not} claim novel cross-domain empirical predictions beyond what each arc supplies.
\item It does \textbf{not} introduce new substrate primitives.
\item ``Cross-domain echo'' = a structural correspondence between the Q-COMPUTE $\mathcal{M}_\mathrm{cap}$ saturation phenomenon and the BH horizon $\Gamma_\mathrm{cross} \to 0$ decoupling phenomenon; both are V5-cross-chain-bandwidth phenomena at substrate level.
\end{enumerate}

\section*{Abstract}

The Q-COMPUTE multiplicity-cap collapse ($\mathcal{M}_\mathrm{cap}(s) \to \mathcal{M}_\mathrm{crit}$, Papers\_053--060) and the BH horizon $\Gamma_\mathrm{cross} \to 0$ decoupling-surface formation (Paper\_039) are FORCED to be structurally echoing phenomena at substrate level. Both arise from V5 cross-chain correlation kernel saturation: $\Gamma_\mathrm{cross}$ collapse marks the substrate-level loss of cross-chain participation flux across the horizon; $\mathcal{M}_\mathrm{cap}$ saturation marks the substrate-level loss of cross-chain correlation budget at the platform's coherence boundary. Both are \textbf{V5-cutoff phenomena} under different platform configurations. The echo is \textbf{structural} (A$\to$position): no Q-COMPUTE platform realizes a horizon, but the substrate-level mechanism is shared.

\section{Primitive Inputs and Paper-Specific Postulates}

\subsection{Primitives invoked (per Paper\_087)}

\textbf{P04 (Bandwidth)}, \textbf{P10 (Rule-type primitive)} [V5], \textbf{P11 (Commitment-irreversibility)}.

\subsection{Upstream dependencies}

\begin{itemize}[noitemsep]
\item \textbf{I-039:} Horizon as decoupling surface ($\Gamma_\mathrm{cross} \to 0$, Paper\_039).
\item \textbf{I-047:} Hawking spectrum (Paper\_047).
\item \textbf{I-050:} Page curve (Paper\_050).
\item \textbf{I-053--060:} Q-COMPUTE arc (Papers\_053--060).
\item \textbf{I-090:} V5 (Paper\_090).
\end{itemize}

\subsection{Paper-specific postulates}

\textbf{P-V5-Shared-Mechanism:} $\Gamma_\mathrm{cross} \to 0$ collapse (BH) and $\mathcal{M}_\mathrm{cap}(s) \to \mathcal{M}_\mathrm{crit}$ collapse (Q-COMPUTE) are governed by the same V5 substrate mechanism --- namely, the cross-chain correlation budget reaching a saturation boundary.

\section*{\S2.5 Audit Table}

\begin{center}
\begin{tabular}{cllp{0.40\textwidth}}
\toprule
\textbf{\#} & \textbf{Step} & \textbf{Label} & \textbf{Notes} \\
\midrule
1 & $\Gamma_\mathrm{cross} \to 0$ at BH horizon & I & Paper\_039. \\
2 & $\mathcal{M}_\mathrm{cap} \to \mathcal{M}_\mathrm{crit}$ in Q-COMPUTE & I & Papers\_053--060. \\
3 & Both phenomena involve V5 saturation & P-V5-Shared-Mechanism & Postulate. \\
4 & BH horizon = substrate cross-chain bandwidth collapse & I & Paper\_039. \\
5 & Q-COMPUTE wall = substrate cross-chain bandwidth saturation & I & Paper\_058. \\
6 & Structural echo (not operational identity) & A$\to$position & Composite. \\
7 & No new cross-domain empirical predictions & I & Each arc supplies own predictions. \\
\bottomrule
\end{tabular}
\end{center}

\section{The Cross-Domain Echo}

\subsection{BH side: $\Gamma_\mathrm{cross} \to 0$}

Paper\_039: the BH horizon is defined as the substrate surface where cross-chain participation flux $\Gamma_\mathrm{cross}$ between inside and outside collapses to zero. This is a V5-saturation phenomenon: V5 cross-chain correlations cannot transmit across the horizon boundary.

\subsection{Q-COMPUTE side: $\mathcal{M}_\mathrm{cap} \to \mathcal{M}_\mathrm{crit}$}

Paper\_058 (Class C plateau): as Q-COMPUTE platforms approach $\mathcal{M}_\mathrm{crit}$, the V5 cross-chain correlation budget saturates. The substrate cannot support additional rule-type configurations without forcing individuation.

\subsection{Same V5 mechanism}

By P-V5-Shared-Mechanism, both are \textbf{the same substrate phenomenon} under different platform configurations:

\begin{itemize}[noitemsep]
\item \textbf{BH:} large-scale gravitational platform; horizon formation = boundary V5 collapse.
\item \textbf{Q-COMPUTE:} laboratory-scale platform; wall formation = budget V5 saturation.
\end{itemize}

The substrate mechanism is V5-cutoff-driven in both cases.

\subsection{What this is NOT}

The echo is \textbf{structural}, not operational. A Q-COMPUTE wall does not produce Hawking radiation; a black hole does not run Shor's algorithm. The shared content is the \textbf{substrate-level mechanism} --- V5 cross-chain saturation forcing individuation / decoupling.

This is an \textbf{A$\to$position} framing claim, not a derivation.

\section{Falsification Criteria}

\begin{itemize}[noitemsep]
\item \textbf{F1:} Substrate evidence that $\Gamma_\mathrm{cross}$ collapse and $\mathcal{M}_\mathrm{cap}$ saturation are governed by \textbf{different} substrate mechanisms --- refutes P-V5-Shared-Mechanism.
\item \textbf{F2:} Empirical detection that V5 finite bandwidth (Paper\_090) is not the load-bearing element in either phenomenon --- refutes the shared-mechanism claim.
\item \textbf{F3:} Discovery of a Q-COMPUTE platform that crosses $\mathcal{M}_\mathrm{crit}$ without V5 saturation --- refutes the substrate identification.
\end{itemize}

\section{Position Statement}

BH $\Gamma_\mathrm{cross}$ collapse and Q-COMPUTE $\mathcal{M}_\mathrm{cap}$ saturation are \textbf{FORCED to share} a substrate-level V5 mechanism under P-V5-Shared-Mechanism. The echo is \textbf{structural} (A$\to$position): same substrate phenomenon, different platforms. Operational identity NOT claimed.

\end{document}
